\documentclass{amsart}

% 数学相关宏包
\usepackage{amsmath,mathtools,amsthm,amssymb}
\usepackage{bbm}
\usepackage{mathrsfs}
\usepackage{centernot}

% 图形和绘图
\usepackage{graphicx,float}
\usepackage{tikz,tikz-cd}
\usepackage{pgfplots}
\usepackage[framemethod=TikZ]{mdframed}
\usetikzlibrary{decorations.pathmorphing,positioning,arrows,automata,calc,shapes,patterns}
\pgfplotsset{compat=1.18}
\usepgfplotslibrary{statistics}

% 表格和列表
\usepackage{booktabs,multirow,colortbl}
\usepackage{enumitem}
\setlist[itemize,enumerate,description]{noitemsep}
\setlistdepth{9}

% 算法
\usepackage{algorithm,algorithmic}

% 字体和颜色
\usepackage{xcolor,listings}
\usepackage{xeCJK}
\setCJKmainfont{Noto Serif CJK SC}

% 页面和链接
\usepackage{fancyhdr,caption}
\usepackage{hyperref}
\usepackage{cleveref}

% 工具宏包
\usepackage{etoolbox,letltxmacro,docmute}
\usepackage{gensymb}

% 自定义设置
\renewcommand{\arraystretch}{1.5}
\let\originalmathbb\mathbb
\renewcommand{\mathbb}[1]{%
    \ifstrequal{#1}{1}{\mathbbm{#1}}{\originalmathbb{#1}}%
}

% 颜色定义
\definecolor{lightblue}{RGB}{173,216,230}
\definecolor{lightgreen}{RGB}{144,238,144}
\definecolor{lightyellow}{RGB}{255,255,224}
\definecolor{lightgray}{RGB}{211,211,211}

% 超链接设置
\hypersetup{
    colorlinks=true,
    linkcolor=red,
    filecolor=magenta,
    urlcolor=cyan,
    pdftitle={\@title},
    pdfauthor={\@author},
    pdfsubject={数学笔记},
    pdfkeywords={数学},
    bookmarksnumbered=true,
    bookmarksopen=true,
    bookmarksopenlevel=6,
    pdfstartview=FitH,
    pdfpagemode=UseOutlines,
    pdfborder={0 0 0}
}

% 定理环境
\theoremstyle{plain}
\newtheorem{theorem}{Theorem}[section]
\newtheorem{lemma}[theorem]{Lemma}
\newtheorem{corollary}[theorem]{Corollary}
\newtheorem{proposition}[theorem]{Proposition}

\theoremstyle{definition}
\newtheorem{definition}[theorem]{Definition}
\newtheorem{example}[theorem]{Example}
\newtheorem{exercise}[theorem]{Exercise}
\newtheorem{problem}[theorem]{Problem}

\theoremstyle{remark}
\newtheorem{remark}[theorem]{Remark}
\newtheorem{note}[theorem]{Note}
\newtheorem{observation}[theorem]{Observation}

% 解答环境
\newtheorem*{solution}{Solution}
\newtheorem*{proof*}{Proof}

% 图片路径
\graphicspath{{F:/iCloudDrive/iCloud~md~obsidian/2025-summer/figures/}}

\author{乐绎华, Yue Yihua}
\address{中山大学, Sun Yat-sen University}
\email{yueyh@mail2.sysu.edu.cn}
\date{\today}
\title{Mathematical Notes}

\begin{document}

\maketitle
\tableofcontents

\section{Grassmannians, Schubert Varieties, Flag Varieties, and Schubert Calculus}

\subsection{Why study Grassmannian?}

\begin{definition}[Grassmannian $\mathrm{Gr}(k, n)$]
The \textbf{Grassmannian} $\mathrm{Gr}(k, n)$ is the set of all $k$-dimensional linear subspaces of an $n$-dimensional vector space $\mathbb{F}^n$ (where $\mathbb{F}$ is usually $\mathbb{R}$ or $\mathbb{C}$).
Formally:
\[
\mathrm{Gr}(k, n) = \{ \text{\(k\)-dimensional subspaces of } \mathbb{F}^n \}.
\]	\begin{itemize}
		\item Example: $\mathrm{Gr}(1, n)$ is the \textbf{projective space} $\mathbb{P}^{n-1}$, since a 1D subspace is just a line through the origin.
		\item Example: $\mathrm{Gr}(n-1, n)$ can be thought of as the set of hyperplanes through the origin.
	\end{itemize}
The Grassmannian is not just a set — it is a \textbf{smooth manifold} (and even an \textbf{algebraic variety} when $\mathbb{F} = \mathbb{C}$).
Its \textbf{dimension} is:
\[
\dim_{\mathbb{F}} \mathrm{Gr}(k, n) = k(n-k).
\]
\end{definition}
Imagine $\mathbb{R}^n$. You can look at:

\begin{itemize}
	\item Points ($k=0$)
	\item Lines through the origin ($k=1$)
	\item Planes through the origin ($k=2$)
	\item etc.
\end{itemize}

The Grassmannian $\mathrm{Gr}(k, n)$ is like a "space whose points are subspaces" — you can \textbf{move smoothly from one subspace to another}. This lets us treat "families of subspaces" as geometric objects themselves.

A key idea:
The Grassmannian is \textbf{homogeneous}:
\[
\mathrm{Gr}(k, n) \cong \frac{\mathrm{GL}(n, \mathbb{F})}{P},
\]
where $P$ is a \textbf{parabolic subgroup} that fixes a chosen $k$-dimensional subspace. This shows that all $k$-planes look "the same" up to a linear transformation.

When $\mathbb{F} = \mathbb{C}$, $\mathrm{Gr}(k, n)$ can be embedded into projective space via \textbf{Plücker coordinates}:
\[
V \in \mathrm{Gr}(k, n) \quad \mapsto \quad v_1 \wedge \cdots \wedge v_k \in \mathbb{P}\big(\Lambda^k \mathbb{C}^n\big),
\]
where $\{v_i\}$ is a basis for $V$.

The image is cut out by \textbf{quadratic Plücker relations}, making the Grassmannian a \textbf{projective algebraic variety}.

This algebraic viewpoint is key in \textbf{Schubert calculus}, where we study how subspaces intersect.

\subsubsection{\textbf{Linear Algebra \& Geometry}}

\begin{itemize}
	\item \textbf{Parametrizing solutions}: If you have a family of subspaces (e.g., solution spaces to a system of equations), they can be described as points in a Grassmannian.
	\item \textbf{Change of basis}: Studying all $k$-subspaces is like studying all possible column spaces of $n \times k$ matrices up to linear transformations.
\end{itemize}

\subsubsection{\textbf{Algebraic Geometry}}

\begin{itemize}
	\item Grassmannians are \textbf{building blocks}: Many moduli spaces can be constructed from them.
	\item \textbf{Schubert varieties} inside Grassmannians help compute intersection numbers (enumerative geometry problems like “how many lines in $\mathbb{P}^3$ meet four given lines?”).
\end{itemize}

\subsubsection{\textbf{Topology \& Characteristic Classes}}

\begin{itemize}
	\item Grassmannians are classifying spaces for vector bundles:
\end{itemize}
\[
\mathrm{Gr}_k(\mathbb{R}^\infty) \quad\text{classifies rank-\(k\) real vector bundles}.
\]
\begin{itemize}
	\item Their cohomology ring is generated by \textbf{Chern/Schubert classes}, essential in topology and geometry.
\end{itemize}

\subsubsection{\textbf{Representation Theory}}

\begin{itemize}
	\item As homogeneous spaces $G/P$, Grassmannians are examples of \textbf{flag varieties}, and their geometry reflects the structure of Lie groups and Lie algebras.
\end{itemize}

\subsubsection{\textbf{Physics}}

\begin{itemize}
	\item In \textbf{quantum field theory} and \textbf{string theory}, Grassmannians appear in:
	\begin{itemize}
		\item \textbf{Gauge theory}: Moduli of certain field configurations are often Grassmannians or their generalizations.
		\item \textbf{Twistor theory}: The "momentum space" in some scattering amplitude formulations is a Grassmannian.
	\end{itemize}
\end{itemize}

\subsubsection{\textbf{Applied Mathematics}}

\begin{itemize}
	\item \textbf{Statistics \& Data Science}: Subspaces model principal components in PCA; the set of all $k$-dimensional subspaces is the search space in \textbf{subspace methods}.
	\item \textbf{Signal Processing}: In array processing and MIMO systems, signals lie in subspaces, and the Grassmannian provides a natural search domain.
\end{itemize}

\begin{table}[h]
	\centering
	\begin{tabular}{|c|c|}
		\hline
		Aspect & Why it matters \\
		\hline
		\textbf{Geometry} & Smooth, compact, homogeneous manifold parametrizing subspaces \\
		\hline
		\textbf{Algebra} & Projective variety with explicit Plücker equations \\
		\hline
		\textbf{Topology} & Classifying space for vector bundles \\
		\hline
		\textbf{Combinatorics} & Schubert calculus counts geometric configurations \\
		\hline
		\textbf{Applications} & From quantum physics to machine learning \\
		\hline
	\end{tabular}
\end{table}
\subsection{Why study Grassmannian?}

\begin{definition}[Grassmannian $\mathrm{Gr}(k, n)$]
The \textbf{Grassmannian} $\mathrm{Gr}(k, n)$ is the set of all $k$-dimensional linear subspaces of an $n$-dimensional vector space $\mathbb{F}^n$ (where $\mathbb{F}$ is usually $\mathbb{R}$ or $\mathbb{C}$).
Formally:
\[
\mathrm{Gr}(k, n) = \{ \text{\(k\)-dimensional subspaces of } \mathbb{F}^n \}.
\]	\begin{itemize}
		\item Example: $\mathrm{Gr}(1, n)$ is the \textbf{projective space} $\mathbb{P}^{n-1}$, since a 1D subspace is just a line through the origin.
		\item Example: $\mathrm{Gr}(n-1, n)$ can be thought of as the set of hyperplanes through the origin.
	\end{itemize}
The Grassmannian is not just a set — it is a \textbf{smooth manifold} (and even an \textbf{algebraic variety} when $\mathbb{F} = \mathbb{C}$).
Its \textbf{dimension} is:
\[
\dim_{\mathbb{F}} \mathrm{Gr}(k, n) = k(n-k).
\]
\end{definition}
Imagine $\mathbb{R}^n$. You can look at:

\begin{itemize}
	\item Points ($k=0$)
	\item Lines through the origin ($k=1$)
	\item Planes through the origin ($k=2$)
	\item etc.
\end{itemize}

The Grassmannian \$\textbackslash{}math...

\subsection{Schubert varieties}

\subsubsection{1. \textbf{What is a Schubert variety?}}

Inside the Grassmannian $\mathrm{Gr}(k,n)$, a \textbf{Schubert variety} is a subset defined by \textbf{linear incidence conditions} with respect to a fixed \textbf{flag} of subspaces
\[
0 = F_0 \subset F_1 \subset F_2 \subset \cdots \subset F_n = \mathbb{F}^n,
\]
where $\dim F_i = i$.

For example, in $\mathrm{Gr}(2,4)$, a Schubert variety might be:

\begin{note}
The set of 2-planes $V$ that intersect a fixed 2-plane $F_2$ in at least a line.
\end{note}
\subsubsection{2. \textbf{Geometric Motivation}}

The real motivation comes from \textbf{enumerative geometry}:

\begin{itemize}
	\item Classical 19th-century geometers asked:
\begin{note}
“How many lines in $\mathbb{P}^3$ meet four given lines in general position?”
\end{note}	\item This problem can be rephrased as:
\begin{note}
“Count the intersection points of certain Schubert varieties in $\mathrm{Gr}(1,4)$.”
\end{note}\end{itemize}

\textbf{Why?}
Each incidence condition (“line meets a given line”) defines a Schubert variety. Intersecting them corresponds to imposing all the conditions at once, and counting intersection points solves the original problem.

\subsubsection{3. \textbf{Why They’re Useful}}

Schubert varieties give us:

\begin{itemize}
	\item \textbf{A combinatorial handle} on the Grassmannian’s geometry:
They are indexed by combinatorial objects (like partitions or Young diagrams), so geometry problems become combinatorics problems.
	\item \textbf{A basis for cohomology/Chow rings}:
The classes $[X_\lambda]$ of Schubert varieties form a basis for $H^*(\mathrm{Gr}(k,n))$. This means \textit{any} intersection problem in the Grassmannian can be reduced to multiplying Schubert classes and reading off coefficients.
	\item \textbf{Natural subvarieties}:
Many geometric constraints are \textit{automatically} Schubert conditions when expressed in a fixed reference flag.
\end{itemize}

\subsubsection{4. \textbf{Broader Motivation}}

\paragraph{1. \textbf{Enumerative Geometry}}

\begin{itemize}
	\item All “counting subspaces with incidence conditions” problems live in the world of Schubert calculus.
	\item Famous example:
	\begin{itemize}
		\item How many 2-planes in $\mathbb{P}^5$ meet 9 general 3-planes?
→ This is the coefficient of the top class in a product of Schubert classes.
	\end{itemize}
\end{itemize}

\paragraph{2. \textbf{Representation Theory}}

\begin{itemize}
	\item Schubert varieties are \textbf{orbit closures} of a Borel subgroup in a flag variety $G/B$.
	\item Their geometry encodes \textbf{weights, Weyl groups, and Bruhat order} — central in Lie theory.
\end{itemize}

\paragraph{3. \textbf{Algebraic Geometry}}

\begin{itemize}
	\item They are typically singular, but with nice \textbf{cell decompositions} (Bruhat cells), which make them a perfect testing ground for singularity theory, intersection theory, and equivariant cohomology.
\end{itemize}

\paragraph{4. \textbf{Combinatorics}}

\begin{itemize}
	\item Multiplication of Schubert classes corresponds to \textbf{Littlewood–Richardson coefficients} — fundamental in symmetric function theory and Young tableaux combinatorics.
\end{itemize}

\paragraph{5. \textbf{Physics}}

\begin{itemize}
	\item Appear in \textbf{string theory} and \textbf{scattering amplitudes}: The positive Grassmannian (a subset where certain Plücker coordinates are positive) is built from Schubert cells.
\end{itemize}

\subsubsection{5. \textbf{Why Study Them? — The Core Reasons}}

\begin{enumerate}
	\item They are the “coordinate charts” for enumerative problems in the Grassmannian.
	\item Their cohomology classes form a canonical basis for intersection theory.
	\item They link geometry, algebra, and combinatorics in a highly structured way.
	\item They are universal examples of orbit closures, hence key in representation theory.
	\item They provide computable, exact answers to problems that would otherwise be messy geometric puzzles.
\end{enumerate}

If you think of the Grassmannian as the \textit{stage}, then \textbf{Schubert varieties are the props} we use to set up very precise geometric conditions. Studying their intersections is like running the “experiments” of enumerative geometry — and Schubert calculus is the toolkit that makes sure we get the exact count every time.

\paragraph{An example of application of Schubert varieties}

\textbf{Problem}: How many lines in projective 3-space $\mathbb{P}^3$ meet four given lines in general position?

\subparagraph{1. Rewriting the problem in Grassmannian language}

A line in $\mathbb{P}^3$ is the same thing as a \textbf{2-dimensional subspace} of $\mathbb{C}^4$ (if we work over $\mathbb{C}$ to avoid degeneracy issues).

So, the space of all lines in $\mathbb{P}^3$ is:
\[
\mathrm{Gr}(1,3) \quad \text{(but in projective indexing that’s usually written as \(\mathrm{Gr}(2,4)\))}
\]
This Grassmannian has complex dimension:
\[
\dim \mathrm{Gr}(2,4) = 2(4 - 2) = 4
\]
\subparagraph{2. The incidence condition}

We fix one line $\ell \subset \mathbb{P}^3$.

\textbf{Condition}: “Line $L$ meets $\ell$” means that:

\begin{itemize}
	\item The corresponding 2-plane $V_L$ in $\mathbb{C}^4$ has \textbf{at least} a 1-dimensional intersection with $V_\ell$ (the 2-plane for $\ell$).
\end{itemize}

This is exactly a \textbf{Schubert variety} in $\mathrm{Gr}(2,4)$:
\[
\Sigma_1 = \{ V \in \mathrm{Gr}(2,4) \mid \dim(V \cap F_2) \ge 1 \}
\]
for a certain 2-plane $F_2$ coming from $\ell$.

\begin{itemize}
	\item Geometrically: It’s the set of all lines in $\mathbb{P}^3$ that meet a fixed line.
	\item Algebraically: $\Sigma_1$ is a \textbf{codimension-1 Schubert variety} in $\mathrm{Gr}(2,4)$.
\end{itemize}

\subparagraph{3. Stacking the conditions}

We now have \textbf{four general lines} $\ell_1, \ell_2, \ell_3, \ell_4$. Each condition “meets $\ell_i$” corresponds to a \textbf{different Schubert variety} $\Sigma_1^{(i)}$ of codimension 1.

The problem reduces to:
\[
\Sigma_1^{(1)} \cap \Sigma_1^{(2)} \cap \Sigma_1^{(3)} \cap \Sigma_1^{(4)} \quad \text{in } \mathrm{Gr}(2,4)
\]
\begin{itemize}
	\item The Grassmannian has \textbf{dimension 4}.
	\item Each $\Sigma_1^{(i)}$ has \textbf{codimension 1}.
	\item Intersecting 4 of them should give us a \textbf{0-dimensional} set — i.e., a finite number of points.
\end{itemize}

\subparagraph{4. Counting via Schubert calculus}

In the cohomology ring $H^*(\mathrm{Gr}(2,4))$, the class $[\Sigma_1]$ is usually denoted $\sigma_1$.

Schubert calculus tells us:
\[
\sigma_1^4 = 2 \cdot [\text{pt}]
\]
where $[\text{pt}]$ is the fundamental class of a point.

This means:
\[
\# \left( \Sigma_1^{(1)} \cap \Sigma_1^{(2)} \cap \Sigma_1^{(3)} \cap \Sigma_1^{(4)} \right) = 2
\]
\subparagraph{5. Interpretation}

We’ve just proven:

\begin{note}
There are \textbf{exactly two} lines in $\mathbb{P}^3$ that meet four general lines in space.
\end{note}
This is a \textbf{nontrivial enumerative fact} — you could try to solve it purely with coordinate geometry, but it’s messy. With Schubert varieties, it’s a clean calculation.

\subparagraph{6. Moral}

This example encapsulates \textbf{why we study Schubert varieties}:

\begin{enumerate}
	\item They \textbf{encode geometric constraints} (like “meets a line”) in a Grassmannian.
	\item Their intersection theory \textbf{counts solutions} to incidence problems.
	\item The combinatorics of Schubert classes gives the count without having to draw messy diagrams or solve big systems of equations.
\end{enumerate}

\subsection{Flag varieties}

\textbf{A flag variety} is the geometric space that parametrizes \textbf{flags} — i.e., nested sequences of subspaces — of a finite-dimensional vector space.

\subsubsection{1. Definition}

Let $V$ be an $n$-dimensional vector space over $\mathbb{F}$ ($\mathbb{R}$ or $\mathbb{C}$).
A \textbf{full flag} is a sequence:
\[
0 \subset V_1 \subset V_2 \subset \cdots \subset V_{n-1} \subset V
\]
where:
\[
\dim V_i = i.
\]
The \textbf{full flag variety} $\mathcal{F}\ell(n)$ is:
\[
\mathcal{F}\ell(n) = \{ \text{all full flags in } V \}.
\]
You can also have \textbf{partial flag varieties} where you only specify certain dimensions:
\[
0 \subset V_{k_1} \subset V_{k_2} \subset \cdots \subset V_{k_r} \subset V, \quad 0 < k_1 < k_2 < \cdots < k_r < n.
\]
Example: The Grassmannian $\mathrm{Gr}(k,n)$ is a partial flag variety with just one subspace.

\subsubsection{2. Group-theoretic viewpoint}

If $G = \mathrm{GL}_n(\mathbb{F})$, then:

\begin{itemize}
	\item A full flag is the orbit of a standard flag under $G$.
	\item The stabilizer of a flag is a \textbf{Borel subgroup} $B$ (upper triangular matrices in the standard basis).
	\item Thus:
\end{itemize}
\[
\mathcal{F}\ell(n) \cong G / B.
\]
More generally, partial flag varieties are:
\[
G / P,
\]
where $P$ is a \textbf{parabolic subgroup}.

\subsubsection{3. Geometry}

\begin{itemize}
	\item The flag variety is a \textbf{smooth projective variety}.
	\item It has a cell decomposition into \textbf{Schubert cells} indexed by permutations (for full flags) or by certain combinatorial data (for partial flags).
	\item These Schubert cells glue into \textbf{Schubert varieties}, which we study for intersection theory.
\end{itemize}

\subsubsection{4. Why It’s Important}

\begin{itemize}
	\item \textbf{Unification}: Many important varieties in algebraic geometry (Grassmannians, quadric hypersurfaces, etc.) are special cases of $G/P$ flag varieties.
	\item \textbf{Representation theory}: Their cohomology and geometry encode information about representations of Lie groups and Lie algebras.
	\item \textbf{Combinatorics}: The Bruhat decomposition of $G/B$ is reflected in the Schubert cell decomposition.
	\item \textbf{Physics}: Appear in gauge theory moduli spaces and in string theory compactifications.
\end{itemize}

\subsubsection{5. Quick Examples}

\begin{itemize}
	\item $\mathrm{Gr}(k,n)$ — one-step partial flag variety.
	\item Full flags in $\mathbb{C}^3$: $0 \subset L \subset P \subset \mathbb{C}^3$
where $L$ is a line, $P$ is a plane containing it.
\end{itemize}

\subsection{What is Schubert calculus?}

Schubert calculus is the \textbf{geometry–algebra dictionary} for answering incidence questions in Grassmannians and more general flag varieties.

\subsubsection{1. \textbf{The problem it solves}}

Imagine a classic enumerative geometry problem:

\begin{note}
How many geometric objects (lines, planes, etc.) satisfy a set of incidence conditions (like “passes through this point” or “meets that line”)?
\end{note}
The \textbf{incidence conditions} are often described by \textbf{Schubert varieties} in a Grassmannian or flag variety.
Schubert calculus gives you an \textbf{algebraic toolkit} to:

\begin{enumerate}
	\item Represent these conditions as elements of the cohomology (or Chow) ring.
	\item Multiply them in that ring.
	\item Read off the number of solutions from the coefficient of the point class.
\end{enumerate}

\subsubsection{2. \textbf{The setup}}

\begin{itemize}
	\item \textbf{Space:} Start with a flag variety $G/P$ — for Grassmannians $\mathrm{Gr}(k,n)$, $G = \mathrm{GL}_n$ and $P$ is a parabolic subgroup.
	\item \textbf{Flags:} Fix a reference flag $F_\bullet$ in $\mathbb{C}^n$.
	\item \textbf{Schubert variety:} For a combinatorial label $\lambda$ (a Young diagram, permutation, or sequence of integers), define:
\end{itemize}
\[
\Omega_\lambda(F_\bullet) \subset \mathrm{Gr}(k,n)
\]
= subvariety of all $k$-planes that meet the steps of $F_\bullet$ in a certain way.

\begin{itemize}
	\item \textbf{Codimension:} $\Omega_\lambda$ has a well-defined codimension $|\lambda|$.
\end{itemize}

\subsubsection{3. \textbf{The algebra side}}

\begin{itemize}
	\item In cohomology $H^*(\mathrm{Gr}(k,n))$, the class $[\Omega_\lambda]$ is denoted $\sigma_\lambda$ and is called a \textbf{Schubert class}.
	\item These $\sigma_\lambda$ form an \textbf{additive basis} for the cohomology ring.
	\item Multiplication:
\end{itemize}
\[
\sigma_\lambda \cdot \sigma_\mu = \sum_{\nu} c_{\lambda,\mu}^\nu \, \sigma_\nu
\]
where $c_{\lambda,\mu}^\nu$ are the \textbf{structure constants} (intersection numbers).

\begin{itemize}
	\item The \textbf{top-degree} class corresponds to a point; if $|\lambda| + |\mu| = \dim\mathrm{Gr}(k,n)$, then:
\end{itemize}
\[
\sigma_\lambda \cdot \sigma_\mu = (\text{integer}) \cdot [\text{pt}]
\]
and that integer is the \textbf{number of solutions}.

\subsubsection{4. \textbf{Computational rules}}

Historically, Schubert calculus came with \textbf{puzzle pieces}:

\begin{itemize}
	\item \textbf{Pieri’s formula}: Multiplication by a special Schubert class corresponding to a single row.
	\item \textbf{Giambelli’s formula}: Express any Schubert class as a determinant in special classes.
	\item \textbf{Littlewood–Richardson rule}: Combinatorial rule for $c_{\lambda,\mu}^\nu$ via tableaux.
\end{itemize}

\subsubsection{5. \textbf{Example}}

From the earlier Grassmannian example ($\mathrm{Gr}(2,4)$):

\begin{itemize}
	\item “Line meets a fixed line” = $\sigma_1$ (codim 1).
	\item Four such conditions: $\sigma_1^4$.
	\item Schubert calculus says $\sigma_1^4 = 2[\text{pt}]$ → there are \textbf{2 lines} meeting four general lines in $\mathbb{P}^3$.
\end{itemize}

\subsubsection{6. \textbf{Why it’s powerful}}

\begin{itemize}
	\item \textbf{Enumerative geometry}: Counts curves, lines, or higher-dimensional planes under constraints.
	\item \textbf{Representation theory}: The cohomology ring structure constants match tensor product multiplicities of $GL_n$ representations.
	\item \textbf{Combinatorics}: Links geometry to Young tableaux, symmetric polynomials, and the symmetric group.
	\item \textbf{Physics}: Appears in gauge theory, quantum cohomology, and mirror symmetry.
\end{itemize}

\subsection{Hodge numbers of Grassmannians}

\subsubsection{Quick summary (the punchline)}

For the complex Grassmannian $\mathrm{Gr}(k,n)$ (a compact complex projective homogeneous manifold):

\begin{itemize}
	\item $H^{p,q}(\mathrm{Gr}(k,n)) = 0$ whenever $p\neq q$.
	\item All nonzero Hodge groups are on the diagonal: $H^{p,p}$.
	\item The Hodge number $h^{p,p} = \dim H^{p,p}(\mathrm{Gr}(k,n))$ equals the number of Schubert classes (equivalently Young diagrams / partitions) of total size $p$ that fit in a $k\times(n-k)$ rectangle.
	\item Equivalently the Poincaré polynomial is the Gaussian binomial (q-binomial) coefficient:
\end{itemize}
\[
P_t(\mathrm{Gr}(k,n))=\sum_{p\ge0} b_{2p} t^{p}=\binom{n}{k}_t,
\]
and $h^{p,p}=b_{2p}$ (the coefficient of $t^p$).

Why this is so and how to actually compute those numbers below.

\subsubsection{Why \texorpdfstring{$H^{p,q}$}{H^p,q} vanishes off the diagonal}

Two quick geometric reasons:

\begin{enumerate}
	\item \textbf{Cell decomposition by Schubert cells.}
$\mathrm{Gr}(k,n)$ admits a Bruhat (Schubert) cell decomposition into complex cells $\mathbb{C}^m$. Every cell has \textit{even real} dimension, so the ordinary cohomology is concentrated in even degrees: $H^{\text{odd}}=0$. For a compact Kähler manifold with only even-degree Betti numbers, the Hodge decomposition forces $H^{p,q}=0$ for $p\ne q$.
	\item \textbf{Homogeneous / Hermitian symmetric structure.}
Grassmannians are projective homogeneous spaces $G/P$; they are Kähler and in fact have Hodge structure of type $(p,p)$ only (a standard fact for rational homogeneous varieties like Grassmannians).
\end{enumerate}

Combining these yields the diagonal-only Hodge numbers.

\subsubsection{How to compute the diagonal Hodge numbers \texorpdfstring{$h^{p,p}$}{h^p,p}}

Here are practical, equivalent approaches.

\paragraph{1 — Count Schubert classes (combinatorial)}

Schubert classes $\sigma_\lambda$ form an additive basis of $H^*(\mathrm{Gr}(k,n))$. They are indexed by partitions (Young diagrams) $\lambda$ whose Ferrers diagram fits inside a $k\times (n-k)$ rectangle. The cohomological degree of $\sigma_\lambda$ is $2|\lambda|$ where $|\lambda|$ is the number of boxes.

Thus
\[
h^{p,p} = \#\{\lambda \mid |\lambda|=p,\ \lambda\text{ fits in }k\times(n-k)\}.
\]
This is the most direct and combinatorial computation.

\paragraph{2 — Use the Gaussian binomial coefficient (Poincaré polynomial)}

The Poincaré polynomial (graded by complex degree) is the q-binomial:
\[
\sum_{p\ge0} h^{p,p} \, t^p \;=\; \binom{n}{k}_t
\;=\; \frac{(1-t^n)(1-t^{n-1})\cdots(1-t^{n-k+1})}{(1-t^k)(1-t^{k-1})\cdots(1-t)}.
\]
Expand this polynomial in $t$; the coefficient of $t^p$ is $h^{p,p}$.

\paragraph{3 — Schur / Chern-class algebra (Giambelli / Schubert calculus)}

Present the cohomology ring in generators (Chern classes of the tautological bundle or quotient bundle) with relations; use Giambelli determinants to express Schubert classes as Schur polynomials. Counting basis elements of degree $2p$ gives $h^{p,p}$.

This is the route you take when you want multiplication rules, not only dimensions.

\paragraph{4 — Representation-theoretic / Borel methods}

Use $G/P$ representation theory:

\begin{itemize}
	\item Bott’s theorem / Borel-Weil-Bott compute cohomology of line bundles.
	\item The decomposition of global sections into irreducible representations explains the grading and yields Betti/Hodge numbers indirectly.
\end{itemize}

This approach is more conceptual and useful when connecting to representation theory.

\paragraph{5 — Computational tools}

For explicit numbers one can:

\begin{itemize}
	\item Enumerate partitions fitting in the box (straightforward programmatically).
	\item Expand the q-binomial coefficient with symbolic algebra (Sage, Mathematica, etc.).
	\item Use known combinatorial formulae (hook-length, etc., in related contexts).
\end{itemize}

\paragraph{Worked example: \texorpdfstring{$\mathrm{Gr}(2,4)$}{Gr(2,4)}}

Here $k=2,n=4$. The rectangle is $2\times 2$. Partitions that fit in $2\times2$ and their sizes:

\begin{itemize}
	\item $\varnothing$ — size 0
	\item $(1)$ — size 1
	\item $(2)$ and $(1,1)$ — size 2 (two shapes)
	\item $(2,1)$ — size 3
	\item $(2,2)$ — size 4
\end{itemize}

Count per $p$:

\begin{itemize}
	\item $p=0:\ h^{0,0}=1$
	\item $p=1:\ h^{1,1}=1$
	\item $p=2:\ h^{2,2}=2$
	\item $p=3:\ h^{3,3}=1$
	\item $p=4:\ h^{4,4}=1$
\end{itemize}

Poincaré polynomial (complex-degree variable $t$):
\[
1 + t + 2t^2 + t^3 + t^4 = \binom{4}{2}_t.
\]
So the only nonzero Hodge numbers are those diagonal entries, matching the combinatorial count.

\paragraph{Practical recommendations}

\begin{itemize}
	\item If you only want the Hodge numbers: use the q-binomial or count partitions fitting in $k\times(n-k)$.
	\item If you want the ring structure (intersection products): learn Schubert calculus (Giambelli, Pieri, Littlewood–Richardson).
	\item If you want Hodge-theoretic or representation-theoretic interpretations: study $G/P$ geometry and Borel–Weil–Bott.
\end{itemize}


\subsection{Cell decomposition}

A \textbf{cell decomposition} (also called a \textbf{CW decomposition} in topology) is a way of building a space by gluing together basic building blocks called \textbf{cells} — spaces homeomorphic to open balls of various dimensions.

In algebraic geometry and Lie theory, when we say a variety has a \textbf{cell decomposition}, we usually mean it can be written as a \textit{disjoint union} of pieces each isomorphic (as a variety) to an affine space $\mathbb{C}^m$.

\subsubsection{1. \textbf{Basic Idea}}

\begin{itemize}
	\item \textbf{0-cells}: points (homeomorphic to $\mathbb{R}^0$ or $\mathbb{C}^0$).
	\item \textbf{1-cells}: open intervals ($\mathbb{R}^1$).
	\item \textbf{2-cells}: open disks ($\mathbb{R}^2$), and so on.
	\item In complex algebraic geometry: we talk about $\mathbb{C}^m$ as the building blocks (even real dimension $2m$).
\end{itemize}

If a variety $X$ has a cell decomposition, it looks like:
\[
X = \bigsqcup_{\alpha} C_\alpha
\]
where each $C_\alpha \cong \mathbb{C}^{d_\alpha}$ is a \textbf{cell}.

\subsubsection{2. \textbf{Why It’s Useful}}

\begin{itemize}
	\item \textbf{Computing topology}:
Each cell contributes one generator to the homology in its dimension; there’s no complicated attaching map in this situation (for smooth projective varieties with such decompositions).
	\item \textbf{Betti numbers}:
If $C_\alpha$ has complex dimension $p$, it contributes to $H^{2p}$ (real degree $2p$), and all odd cohomology vanishes.
	\item \textbf{Intersection theory}:
For homogeneous spaces like Grassmannians, the cells correspond to \textbf{Schubert cells}, and their closures are \textbf{Schubert varieties} — the natural generators of the cohomology ring.
\end{itemize}

\subsubsection{3. \textbf{Example: Projective space}}

$\mathbb{CP}^n$ has a cell decomposition:
\[
\mathbb{CP}^n = \mathbb{C}^0 \;\sqcup\; \mathbb{C}^1 \;\sqcup\; \mathbb{C}^2 \;\sqcup\; \cdots \;\sqcup\; \mathbb{C}^n
\]
given by:
\[
[z_0: \dots : z_n], \quad \text{cell of dim }m = \{ [1: z_1 : \dots : z_m : 0 : \dots : 0] \}
\]
So $H^{p,q}=0$ unless $p=q$, and $h^{p,p}=1$ for $0 \le p \le n$.

\subsubsection{4. \textbf{Example: Grassmannian \texorpdfstring{$\mathrm{Gr}(2,4)$}{Gr(2,4)}}}

The \textbf{Schubert cells} are:

\begin{itemize}
	\item $\mathbb{C}^0$ (dimension 0)
	\item $\mathbb{C}^1$ (dim 1)
	\item $\mathbb{C}^2$ (two distinct ones)
	\item $\mathbb{C}^3$
	\item $\mathbb{C}^4$
\end{itemize}

These cells glue together to make $\mathrm{Gr}(2,4)$, and each closure is a Schubert variety.

\subsubsection{5. \textbf{Connection to Schubert calculus}}

In $G/P$ (like Grassmannians), the \textbf{Bruhat decomposition}:
\[
G/P = \bigsqcup_{w \in W^P} BwP/P
\]
gives a decomposition into \textbf{Schubert cells}, each isomorphic to $\mathbb{C}^m$. Counting these by dimension directly gives the Betti numbers.


\end{document}